
\begin{keybox}[The sheet generator $V$ is a binary internal compiler \mE\ (\veri{v410\_sheet\_generator\_binary.py})]
The sheet axis $V=Q\operatorname{diag}(0,1,1)$ is not only the rank-$2$ direction of the cross --- its
\emph{iterates} print the binary carrier spine. The closed form (proved by the inductive step
$F(n)V=F(n{+}1)$)
\[
  V^n=\begin{pmatrix}0&2^{n-1}&0\\0&2^{n}&0\\0&2^{n+1}-2&1\end{pmatrix},\qquad
  V^n\mathbf 1=(2^{n-1},\,2^{n},\,2^{n+1}-1),
\]
gives $V^{1..4}\mathbf 1=(1,2,3),(2,4,7),(4,8,15),(8,16,31)$ --- the $A_3$ exponents, then
$(|\Z_2|,|\mu_4|,7)$, $(|\mu_4|,\operatorname{rank}E_8,\dim A_3)$ and
$(\operatorname{rank}E_8,\dim S^+,2^{\gcar}{-}1{=}248/8)$. The four bilinear families are exact in $n$,
\[
  \mathbf 1^TV^n\mathbf 1=7\cdot2^{n-1}-1,\quad \mathbf 1^TV^na=7\cdot2^{n-1},\quad
  a^TV^na=11\cdot2^{n-1},\quad a^TV^n\mathbf 1=11\cdot2^{n-1}-2,
\]
so the recurring flavor integers are one operator's readouts: $(6,7,9,11)$ at $n{=}1$, then
$\mathbf 1^TV^2\mathbf 1=13=\Delta_Q$, $\mathbf 1^TV^3\mathbf 1=27=\mathbf 1^TRa$,
$\mathbf 1^TV^4\mathbf 1=55$ (the quark numerator) and $\mathbf 1^TV^4a=56=\dim\mathbf{56}_{E_7}$; the
amplitude $a^TVa=11$ sits one level below the quark Pl\"ucker norm. \emph{Typing:} the closed form and
the bilinears are exact \mE; the binary spine is \emph{forced} by $\operatorname{Spec}(V)=\{0,1,2\}$ (the
dominant eigenvalue $2$ doubles, the eigenvalue $1$ fixes the third slot), so the carrier/Lie matches are
audit-typed \mC, exactly like the $G_2$ center reading above.
\smallskip\\
\noindent\textbf{The quark $u/d$ ratio is a pure $V$-power readout (\veri{v411\_ud\_ratio\_vpower.py}).}
The closed quark ratio re-expresses as a quotient of three bilinear readouts of the \emph{one} generator,
\[
  \boxed{\ \frac{c_u}{c_d}=\frac{\mathbf 1^TV^4\mathbf 1}{(a^TV\mathbf 1)(\mathbf 1^TV^2\mathbf 1)}
  =\frac{55}{9\cdot13}=\frac{55}{117}\ },
\]
identical to the established $\gcar\|\operatorname{Pl}(K)\|_1/(\Nfam^2\Delta_Q)=5\cdot11/(9\cdot13)$ --- the
numerator merely repackages $\gcar\cdot11=55=\mathbf 1^TV^4\mathbf 1$. This is an exact \emph{re-encoding}:
the value and the $u/d$ identification stay \mC, coupled to the $(U_{\mathrm{wall}})$ readout rigidity.
\smallskip\\
\noindent\textbf{The source corner $J$ on the $\Z_2$ wall (\veri{v412\_sheet\_source\_corner\_J.py}).} The
operator on the $t{=}{-}2$ iso-volume wall (the $\Z_2$ layer $\det{=}2$, \vref{v135}), so far unnamed,
$J:=M(1,{-}2)=C-2V=\big(\begin{smallmatrix}4&1&0\\4&1&2\\5&1&1\end{smallmatrix}\big)$, completes the sheet
axis $J,K,C,F=M(1,{-}2..1)$ and is tightly typed:
\[
  \chi_J=\lambda^3-6\lambda^2+3\lambda-2\ \Rightarrow\ (\operatorname{tr},e_2,\det)=(6,3,2)
  =(|R^+(A_3)|,\Nfam,|\Z_2|),
\]
with $a^TJa=30=h(E_8)$, $\det(I{+}J)=12=\dim\mathfrak g_{\mathrm{SM}}=\chi_{E_8}(i)$ and
$\det(2I{+}J)=40=|R(D_5)|$ --- the Lie readings audit-typed like the $G_2/F_4$ shadows.
\smallskip\\
\noindent\textbf{The sheet axis is a characteristic-difference calculator
(\veri{v413\_sheet\_characteristic\_calculus.py}).} On $M_t=C+tV$ the elementary-symmetric coefficients
encode the atoms as \emph{difference orders}: $e_1(t)=3t{+}12$, $e_2(t)=2t^2{+}15t{+}25$,
$e_3(t)=\det M_t=4t^2{+}14t{+}14$, so $\Delta e_1=3=\Nfam$, $\Delta^2 e_2=4=|\mu_4|$ and
$\Delta^2 e_3=8=\operatorname{rank}E_8$ (the $e_3$ curvature is the \vref{v218}/\vref{v224} sheet
quadratic; the full triple is new). The anchor energy is exactly linear with the amplitude $11=a^TVa$,
\[
  a^TM_t a=52+11t:\qquad 30\xrightarrow{+11}41\xrightarrow{+11}52\xrightarrow{+11}63
  =(h(E_8),\,10b_1,\,\dim F_4,\,7\Nfam^2),
\]
so $11$ is the \emph{step} of the anchor energy along the sheet axis, not only the quark Pl\"ucker norm.
\smallskip\\
\noindent\textbf{The center is a resolvent portal $G_2\!\to\!F_4\!\to\!E_8$
(\veri{v414\_center\_resolvent\_portal.py}).} The center $C=M(1,0)$ (the $t{=}0$ value of $V$) has
$\operatorname{tr}C=12=\dim\mathfrak g_{\mathrm{SM}}=\chi_{E_8}(i)$, $\det C=14=\dim G_2$ and
$\sum_{ij}C_{ij}=31=2^{\gcar}{-}1=248/8$; its shifted determinants climb the exceptional shell ladder
\[
  \det C=14=\dim G_2,\qquad \det(I{+}C)=52=\dim F_4,\qquad \det(2I{+}C)=120=|R^+(E_8)|,
\]
i.e.\ $\det(C{+}kI)$ for $k=0,1,2$ traces $G_2\!\to\!F_4\!\to\!E_8^+$ --- exact \mE, the portal reading \mC.
\smallskip\\
\noindent\textbf{The center quotient compiler: the atoms as one operator, and the CP-halved Coxeter code
(\veri{v530\_center\_quotient\_compiler.py}).} The center carries a \emph{fixed primitive line}: with
$v=(-1,1,0)$ one has $Cv=v$ exactly, and in the unimodular basis $P=(v,e_1,e_3)$ the induced operator on
the quotient lattice $\Z^3/\Z v$ is
\[
  A=\begin{pmatrix}8&2\\5&3\end{pmatrix}
   =\begin{pmatrix}\operatorname{rank}E_8&|\Z_2|\\ \gcar&\Nfam\end{pmatrix}:
\]
the four discrete atoms form \emph{one integer quotient operator}, with
$\operatorname{tr}A=11$, $\det A=14=\det C$, $\operatorname{disc}\chi_A=65=5\cdot13$,
$\operatorname{perm}A=34=p_5(a)$, $\prod A_{ij}=240=|R(E_8)|$ and $\det(A{-}I)=4=|\mu_4|$. The center
\emph{self-codes}: its row sums $C\mathbf 1=(7,11,13)$ equal exactly
$(\det A/|\Z_2|,\ \operatorname{tr}A,\ \operatorname{disc}\chi_A/\gcar)$, and together with the fixed
eigenvalue $1$ this is $(1,7,11,13)$ --- the positive representatives of the CP pairs $m\sim30{-}m$ of the
$E_8$ Coxeter exponents $(\Z/30)^\times$ (the CP-halved Coxeter clock, \vref{v223}). The \vref{v414}
resolvent ladder acquires its inner structure: $\det(C{+}kI)=(k{+}1)\det(A{+}kI)$ \emph{identically}, with
$(k{+}1)=(1,2,3,4)=(N_\Phi,|\Z_2|,\Nfam,|\mu_4|)$ and the quotient ladder
$\det(A{+}kI)=(14,26,40,56)$ (readings $\dim G_2$, $\dim V_{26}(F_4)$, $|R(D_5)|$, $\dim V_{56}(E_7)$,
audit-typed \mC), first differences $(12,14,16)=(\dim\mathfrak g_{\mathrm{SM}},\dim G_2,\dim S^+)$, and the
new endpoint $\det(C{+}3I)=224=248-24$. \emph{Honest uniqueness:} eigenvalue $1$ is \emph{not} unique to
$C$ ($K$ and $Q$ share it), but the $K$ quotient
$\big(\begin{smallmatrix}14&8\\-11&-6\end{smallmatrix}\big)$ and the $Q$ quotient
$\big(\begin{smallmatrix}3&1\\3&2\end{smallmatrix}\big)$ are not the atom matrix and cannot self-code
($5\nmid48$; $3/2\notin\Z$) --- $C$ is unique in the \emph{conjunction} fixed line $+$ atom quotient $+$
CP self-code; the full selector-stratum ablation stays open \mO. Exact \mE, readings \mC.
\emph{Bridge corollary \mE\ (\veri{v977\_transfer\_unistochastic\_wilson.py}):} the discriminant carries a
Gaussian square-sum structure the sibling quotients lack,
$\operatorname{disc}\chi_A=65=\det(A{-}I)^2+(\det A/2)^2=4^2+7^2$ with $-4+7i=(1{+}2i)(2{+}3i)$
(norms $5\cdot13$), and $e^{i\delta}=(-\det(A{-}I)\pm i\det A/2)/\sqrt{\operatorname{disc}\chi_A}
=(-4\pm7i)/\sqrt{65}$ is \emph{exactly} the phase of the unistochastic $SU(3)$ lift of the seam
transport root $B$ (Jarlskog $J=\pm1/27$; the $K$/$Q$ sibling quotients fail the square-sum identity,
$48\neq3^2{+}2^2$ and $13\neq1^2{+}(3/2)^2$) --- an exact transport--center bridge; the intertwiner that
would force the lift \emph{angles} from $A$-data stays open (research-contracts document,
\texttt{TRANSFER.COHERENT.WILSON.01}).
\smallskip\\
\noindent\textbf{The sheet transfer is a norm line of discriminant $-7$
(\veri{v533\_ftransfer\_disc7\_norm.py}).} The branch-preserving transfer path $J,K,C,F=M(1,-2..1)$
(\vref{v224}/\vref{v412}) carries the exact norm identity
\[
  D(t)=\det M(1,t)=4t^2+14t+14=\frac{(4t+7)^2+7}{4}
      =N_{\mathbb Q(\sqrt{-7})/\mathbb Q}(\alpha_t),\qquad \alpha_t=\frac{4t+7+\sqrt{-7}}2,
\]
with $\alpha_t$ an \emph{algebraic integer} ($-7\equiv1\bmod4$, $4t{+}7$ odd, so
$\alpha_t\in\Z[(1{+}\sqrt{-7})/2]$) and the whole path one arithmetic \emph{translation}
$\alpha_{t+1}=\alpha_t+2=\alpha_t+|\Z_2|$ with norms $(2,4,14,32)=(\det J,\det K,\det C,\det F)$. The
characteristic polynomial $x^2-(4t{+}7)x+D(t)$ has discriminant $-7$ for \emph{every} $t$; the vertex sits
at $t=-7/4$ with $D_{\min}=7/4>0$ (no real zero on the branch-preserving path), and the two sheet
coefficients reappear arithmetically: curvature $D''=8=\operatorname{rank}E_8$, $\sqrt{-\operatorname{disc}}
=\sqrt7$ (the scalaron as fixed quadratic discriminant, \mC). Since $-7\equiv1\bmod 8$, the prime $2$
\emph{splits} in $\mathbb Q(\sqrt{-7})$ (class number $1$): the $J$ endpoint norm $2=|\Z_2|$ is the split prime
above~$2$ (audit). Negative controls: the winding axis $\det M(s,0)=6s{+}8$ is linear and the anchor block
$\det B(1,t)=3t^2{+}21t{+}28$ has \emph{positive} discriminant $105$ --- only the sheet determinant is a
complex norm line. Whether the four external transfer solvers act on $\alpha_t$ by a \emph{common}
fractional-linear action is the preregistered next test (kill: four incompatible arithmetic actions);
$F_{\mathrm{transfer}}$ itself stays an external interface \mO. Exact \mE, the arithmetic frame \mC/\mO.
\smallskip\\
\noindent\textbf{The square CM-norm dictionary is operator-realised, and the deck is intrinsic
(\veri{v415\_gaussian\_operator.py}).} The non-commutativity of the two axes is itself a $\mu_4$
quarter-turn: the integer commutator
\[
  J:=\tfrac13[U,V]=\begin{pmatrix}-1&1&0\\-2&1&0\\-3&1&0\end{pmatrix},\qquad
  \operatorname{Spec}(J)=\{i,-i,0\},\quad J^2=-I\ \text{on the rank-2 image},
\]
is a complex structure born from the binary $V$ and ternary $U$ compilers. It lifts the Gaussian
arithmetic of \vref{v222}/\vref{v230} from numbers to operators: $3I{+}2J$ has spectrum
$\{3{+}2i,3{-}2i,3\}$ and $5I{+}4J$ has $\{5{+}4i,5{-}4i,5\}$, so the Gaussian integers $3{+}2i$ and
$5{+}4i$ are literal eigenvalues, and
\[
  (3I{+}2J)(3I{-}2J)=9I-4J^2\sim\{13,13,9\},\qquad (5I{+}4J)(5I{-}4J)=25I-16J^2\sim\{41,41,25\},
\]
i.e.\ $\Delta_Q=13$ and $10b_1=41$ on the $\mu_4$-plane, the real part$^2$ ($\Nfam^2$, $\gcar^2$) on the
kernel. Finally $D:=-I+J-J^2$ is an \emph{intrinsic} order-$4$ deck ($D^4=I$, $\chi_D=(x{+}1)(x^2{+}1)$,
$\operatorname{Spec}\{i,-1,-i\}$) whose self-conjugate $\chi_2$-line is canonically
$\ker[U,V]=(0,0,1)=a-\mathbf 1$; this exhibits the $\mu_4$ deck as a diamond object and sharpens the
\vref{v140} $\Z_3$ residue to one concrete identification --- the operator identities \mE, the
geometric-deck match \mC.
\smallskip\\
\noindent\textbf{The Eisenstein dual and the CP clock (\veri{v417\_eisenstein\_cp\_operator.py}).}
The same lift works on the \emph{hexagonal} side. The family rotation $P=(1\,2\,3)$ is the order-$3$
Eisenstein deck ($P^2+P+I=\mathrm{ONES}$, so $\omega^2+\omega+1=0$ on the plane $\perp(1,1,1)$ --- the
order-$3$ dual of $J^2=-I$), and it realises the hexagonal CM norm as an operator,
\[
  (3I+2P)(3I+2P^2)=7I+6\,\mathrm{ONES}\sim\{7,7,25\},
\]
i.e.\ $7=N_{\Z[\omega]}(3{+}2\omega)=$ scalaron on the $\omega$-plane and $25=\gcar^2$ on $(1,1,1)$ ---
the exact dual of the square $\{13,13,9\}$. The CP phase is the order-$6$ clock $W=-P^2$ (the triality
$\rho=\omega^2{\cdot}({-}1)$): $W^2=P$, $W^3=-I$, $\operatorname{Spec}W=\{-1,\zeta_6,\zeta_6^{-1}\}$, and the
two measured phases are arguments of its powers, $\delta_{\mathrm{CKM,lead}}=\arg\zeta_6=\pi/3$ and
$\delta_{\mathrm{PMNS}}=\arg\zeta_6^4=4\pi/3$ (\vref{v316}/\vref{v320}). All four flavor-side clocks are
powers of one Coxeter unit $\zeta_{30}$ ($h(E_8)=30=|\Z_2|\Nfam\gcar$): $\zeta_{30}^{15}=-1$,
$\zeta_{30}^{10}=\omega$, $\zeta_{30}^{6}=\zeta_5$, $\zeta_{30}^{5}=\zeta_6$, while the seam $\mu_4=i$ is
\emph{not} a power ($4\nmid30$) but the Galois side $(\Z/30)^\times$ of order
$\varphi(30)=8=\operatorname{rank}E_8$ --- so $\mathbb Q(\zeta_{30})$ carries both $E_8$ invariants at
once. \emph{Honest:} the carrier $\mu_5$ has no $3{\times}3$ operator (the golden $\varphi=2\cos\tfrac\pi5$
is an output-side trace, \vref{v313}/\vref{v349}); the operator/arithmetic identities are \mE, the
clock-assignment readings \mC.
\end{keybox}

\begin{keybox}[The $Q_+$ grading is cohomology: one spectrum in four readouts \mE\ (\veri{v137\_qplus\_cohomology.py})]
The sheet axis is not just numerology on $V$: it is the cohomology of the seam curve.
$b_1(\PP\setminus\mu_4)=3=\Nfam$, with holomorphic basis $\omega_k=z^{k-1}dz/(z^4-1)$, $k=1,2,3$
(regular at infinity; the trivial character is absent), and under the $\mu_4$ deck $z\mapsto iz$
\[
  \omega_k\ \longmapsto\ i^k\,\omega_k,\qquad
  \operatorname{res}_{\zeta}\omega_k=\tfrac{\zeta^k}4 :
\]
$H^1=\chi_1\oplus\chi_2\oplus\chi_3$ --- each nontrivial $\mu_4$ character exactly once, and the residue
vector over the punctures \emph{is} the character vector. The character degrees $(1,2,3)$ are exactly
$\operatorname{Spec}(Q_+)$ ($Q_+=3\,\mathrm{cusp}+1$, \vref{v69}/\vref{v72}), and the cusp weights $\{0,\tfrac13,
\tfrac23\}$ underneath are $\operatorname{spec}(A_0^\star)$ (\vref{v115}):
\[
  \boxed{\ \text{MS weights}\ =\ \text{cusp classes}\ =\ H^1\text{ characters}\ =\
  \operatorname{Spec}(Q_+)-1\ \text{over } \Nfam\ }
\]
--- one spectrum in four readouts; and $(1,2,3)$ are the \emph{exponents of $A_3$} ($W(A_3)=S_4$, \vref{v117};
$h(A_3)=4=|\mu_4|$): the $Q_+$ generation grading \emph{is} the $A_3$ exponent grading, realised on the
seam curve's cohomology. \emph{Typing:} all spectrum statements \mE; what remains of \texttt{GATE.QGEO}
is the \emph{naturality} of the identification --- the canonical map from the $H^1$ character basis to
the $V$-axis generation basis ($\operatorname{rowsum}(V)=(1,2,3)$ verbatim, \veri{v135\_det\_surface.py})
--- typed \mC.
\smallskip\\
\noindent\textbf{The canonical map exists and is rigid (\veri{v140\_canonical\_map.py}).} The naturality
residue collapses. Of the established $\mu_4$ actions on generation space, the plain deck
$U=\operatorname{diag}(1,i,-i)$ has characters $(0,1,3)$ --- an invariant line, missing $2$ --- and
\emph{cannot} match $H^1$; but the \textbf{sheet-twisted} deck $-U$ (characters $(2,3,1)$) and the cusp
exponential $i^{Q_+}$ (characters $(1,2,3)$) both carry the character \emph{set} $\{1,2,3\}$ --- exactly
the $H^1$ set, by the same sign twist that builds the \vref{v118} hexagon. Both sides are multiplicity-free, so
by Schur the equivariant isomorphism \emph{exists and is unique} up to one scalar per character: the
canonical map is well-defined --- existence and uniqueness are no longer open. Under the $i^{Q_+}$-map
the $H^1$ Euler degree pulls back to $Q_+$ itself; under the $-U$-map to $Q_+\!\circ\rho$ with $\rho$ the
$(123)$ family rotation: the two candidates differ by \emph{exactly one} $\Z_3$ rotation. \emph{Typing:}
character algebra, rigidity and grading transport \mE.
\smallskip\\
\noindent\textbf{The deck selection theorem: the $\Z_3$ choice is \emph{derived}
(\veri{v141\_deck\_selection.py}).} The discrete choice does not survive contact with the integer model
the suite already owns: the derived integer deck $G=T_A\Sigma$ (\veri{v97\_sheet\_conjugation\_bridge.py}/%
\veri{v98\_discriminant\_dictionary.py}) has the $Q_+{=}1$ cusp line as an \emph{exact} eigenline with
eigenvalue $-1=i^2$ --- the self-conjugate character $2$ --- and the conjugate pair $\{1,3\}$ on the
$E$-plane. Of the three $\Z_3$ rotations only the sheet-twisted assignment $(2,3,1)=\mathrm{chars}(-U)$
puts character $2$ on the $Q_+{=}1$ line: \textbf{the geometric boundary deck is the sheet-twisted class;
the cusp exponential $i^{Q_+}$ is excluded} (same eigenvalue multiset $\{-1,\pm i\}$, wrong grading
pairing --- the $(-1)$-eigenlines sit on $Q_+$ degrees $1$ vs $2$). An explicit $Q_+{=}1$-line-preserving
conjugacy $SGS^{-1}=-U$ exists; under the relabeling $k\to-k$ only the plane orientation flips $=$ the
established sheet $\Z_2$ (the two Lagrangian glues, \vref{v92}) --- no new residual. Consequence: the
$H^1$ Euler degree pulls back to $Q_+\!\circ\rho$; cusp grading and Euler grading differ on generation
space by exactly the one family rotation $\rho=$ the $\operatorname{coker}Q$ triality (\vref{v72}).
\emph{Typing:} selection, conjugacy and obstruction \mE\ (lattice-level theorem); what remains of
\texttt{GATE.QGEO} is only the realisation premise the gate always carried (the integer $D_4$ model
\emph{is} the parabolic geometry, \vref{v69}/\vref{v98}) \mC\ --- with \emph{no} discrete freedom left.
\smallskip\\
\noindent\textbf{The M\"obius $D_4$ realisation: the premise at its floor
(\veri{v146\_moebius\_d4.py}).} The realisation premise itself now has every checkable part checked.
The \emph{full} dihedral symmetry group of the seam curve, $\langle\delta\!:z\mapsto iz,\ \iota\!:z\mapsto
1/z\rangle=D_4$, acts on the explicit cohomology basis $\omega_k=z^{k-1}dz/(z^4{-}1)$ exactly as the
integer model: $\delta^*\omega_k=i^k\omega_k$ (the deck, \vref{v137}), and
\[
  \boxed{\ \iota^*\omega_1=\omega_3,\quad \iota^*\omega_2=\omega_2,\quad \iota^*\omega_3=\omega_1\ }
\]
--- the exponent duality with parity $+1$ on the self-conjugate line and $\det=-1$, \emph{no extra
phases}. In the cusp basis the anchor-forced conjugation $T_A$ \emph{is} this permutation (fixes the
character-$2$ line with $+1$, $\det=-1$ $=$ the glue-swap parity), and $\Sigma$ \emph{is} the
$\delta\iota$ class ($-1$ on the character-$2$ line, $\det=+1$ $=$ glue-fix): the \vref{v98}
reflection-class split lands exactly on the geometric reflection classes. \emph{Typing:} the dihedral
match \mE; with spectra (\vref{v137}), cyclic part (\vref{v141}) and reflections now all matched, the
premise reduces to the module identification itself --- ``generation space \emph{is} $H^1$ of the seam
curve with its parabolic weights'' \mC\ --- with no remaining checkable mismatch.
\end{keybox}

\begin{winbox}[Spectra and the determinant ladder $\to|W(D_5)|$ \mE]
The two operators carry the compiler atoms in their characteristic polynomials,
\begin{equation}
  \chi_Q=(t-1)(t^2-5t+3)\ \ [\,\gcar,\Nfam\,],\qquad
  \chi_K=(t-1)(t^2-8t+4)\ \ [\,h(D_5),|\mu_4|\,],
\end{equation}
with discriminants $\Delta_Q=13=|R(A_3)|{+}1$ and $\Delta_K=48=\Oadm$. Their determinants form a torsion
ladder ($\operatorname{coker}=\Z_3,\Z_4,\Z_8,\Z_{20}$),
\begin{equation}
\begin{gathered}
  (\det Q,\det K,\det R,\det L)=(3,4,8,20),\\[2pt]
  \boxed{\ \det Q\cdot\det K\cdot\det R\cdot\det L=1920=|W(D_5)|\ }\qquad\mE ,
\end{gathered}
\end{equation}
tying flavor residues, mass powers, transport and the $D_5$ Weyl group (the same $1920$ that appears in the
horizon Hawking-power denominator). The sheet decomposition is sharp: $\operatorname{Spec}(Q_+)=\{1,2,3\}$
are the $A_3$ exponents ($\det Q_+=6=|R^+(A_3)|$), $\chi_{Q_-}=t(t^2-3)$ (a square root of $\Nfam$), and
$\chi_{K_-}=t(t^2-8)$ (a square root of $h(D_5)$). Finally the EM budget sits in the $K$ anchor block:
$a^\top K a=41=10b_1$ with $a=(1,1,2)$, and $\det\!\begin{psmallmatrix}\mathbf1^\top K\mathbf1&\mathbf1^\top Ka\\ a^\top K\mathbf1&a^\top Ka\end{psmallmatrix}=\det\begin{psmallmatrix}25&29\\35&41\end{psmallmatrix}=10=\AL$.
\veri{v10\_projection\_involution.py}, \veri{v12\_mass\_generation\_polynomials.py}
\smallskip\\
\noindent\textbf{Mass volume and the EM budget are \emph{one} anchor operation
(\veri{v74\_compiler\_micro\_lemmas.py}).} The two diagonal entries of $B_K$ are the two physical readouts:
$\mathbf1^\top K\mathbf1=25=\gcar^2$ is the grand \emph{mass-volume} exponent ($\det M_{\rm SM}\!\sim\!(\phiz)^{25}$),
and $a^\top K a=41=10b_1$ is the \emph{EM budget}. Their difference is one generation,
\[
  \boxed{\ a^\top K a-\mathbf1^\top K\mathbf1=41-25=16=\dim S^+\ },
\]
i.e.\ \emph{EM budget $=$ mass volume $+$ one generation}. So the number $41$ that the EM fixed point reads is
not ``$40{+}1$'' but the anchor quadratic form of the mass-power matrix, tied to the mass volume $\gcar^2$ by
a single half-spinor $\dim S^+$.
\end{winbox}

\begin{winbox}[Mass and transport are one pencil: $\det(K+xQ)$ carries the scalaron sequence \mE]
Because $L=K+Q$, the mass-power matrix $K$ and the transport matrix $L$ are the two endpoints
($x=0,1$) of a \emph{single} one-parameter compiler pencil $P(x)=K+xQ$, whose determinant is
\[
  \boxed{\ \det(K+xQ)=3x^3+7x^2+6x+4\ },\qquad
  (3,7,6,4)=\bigl(\Nfam,\ \Oadm-10b_1,\ |R^+(A_3)|,\ |\mu_4|\bigr),
\]
i.e.\ the coefficient sequence is $N_{\mathrm{fam}}$, the \emph{scalaron exponent} $7$, the $A_3$ positive
roots and the glue index --- one curve linking flavor, the Starobinsky scalaron, the family and the glue.
The pencil is anchored at four integer nodes by compiler atoms:
$P(-1)=2=|\Z_2|$ (the $K{-}Q$ sheet endpoint, with $\det(K{-}Q)=|\Z_2|$, $\operatorname{tr}(K{-}Q)=\Nfam$),
$P(0)=4=\det K=|\mu_4|$, $P(1)=20=\det L$, and $P(2)=68=2\,p_5(a)$ (twice the anchor power sum $p_5=34$).
\veri{v52\_pencil\_endpoints.py}
\smallskip\\
\noindent\textbf{And the consecutive \emph{differences} are the sheet$\to$matter chain
(\veri{v74\_compiler\_micro\_lemmas.py}).} The node values $(2,4,20,68)$ step by
\[
  \boxed{\ P(0){-}P(-1)=2=|\Z_2|,\quad P(1){-}P(0)=16=\dim S^+,\quad P(2){-}P(1)=48=\Oadm\ }
\]
i.e.\ $2\to16\to48$ --- \emph{sheet $\to$ one generation $\to$ three generations} --- so the single
determinant pencil already encodes the passage from the $\Z_2$ sheet to the one-family half-spinor to the
full admissible matter packet. The $(K,Q,L)$ operators are therefore not separate artefacts but one curve.
The non-commutativity of the mass$\to$transport step is a fixed $G_2$-sized obstruction:
$[K,L]=[L,Q]=[K,Q]$ (forced by $L=K+Q$) with $\chi_{[K,Q]}(t)=t^3+5t+14$, reading $\gcar=5$ and
$\dim G_2=14$ (audit-level \mO). \veri{v37\_plucker\_anchor.py}
\end{winbox}

\begin{keybox}[Operator-pencil geometry: the anchor plane as a small area in operator space \mE\ (\veri{v80\_operator\_pencil\_geometry.py})]
Read the four operators through the \emph{anchor block} $B_M=\begin{psmallmatrix}\mathbf1^\top M\mathbf1&\mathbf1^\top Ma\\a^\top M\mathbf1&a^\top Ma\end{psmallmatrix}$.
\textbf{(i) Anchor Singularity Theorem.} On the pencil $P(x)=K+xQ$,
\[
  \boxed{\ \det B_{K+xQ}=9x^2+21x+10=(\Nfam x+|\Z_2|)(\Nfam x+\gcar)=(3x+2)(3x+5)\ }\qquad\mE ,
\]
so the oriented anchor area \emph{degenerates} at exactly $x=-|\Z_2|/\Nfam=-\tfrac23$ and $x=-\gcar/\Nfam=-\tfrac53$.
The first zero is the Koide target / transport-gap base ($\lambda_2=(\tfrac23)^6$); the second is the $D_5/A_3$
glue asymmetry. \textbf{So $\tfrac23$ is not just a recurring fraction --- it is the first singular point of the
mass$\leftrightarrow$transport anchor plane}, a structural \emph{location} where an oriented area collapses.
\textbf{(ii) Type checker.} The anchor-block determinants classify the operators:
$\det B_Q=9=\Nfam^2$ (family square), $\det B_K=10=\AL$ (pair sector), $\det B_R=16=\dim S^+$ (one generation),
$\det B_L=40=|R(D_5)|$ (carrier root budget) --- each a determinant of a fixed block, not a loose integer.
\textbf{(iii) Buildup chain.} The pencil coefficients $(3,7,6,4)$ have partial sums $(3,10,16,20)=(\Nfam,\AL,\dim S^+,\det L)$,
so the scalaron exponent is the internal flavor transition $7=\AL-\Nfam$.
\smallskip\\
\noindent\emph{Audit layer (exact values \mE, decorative reading \mC).} $D(x)=\det(K{+}xQ)$ and its derivatives
at $x=-1,0,1,2$ give $(2,4,20,68),(1,6,29,70),(-4,14,32,50)$ ($29$ the top $E_8$ exponent, $32=2^{\gcar}$, $50$
the $A_4{\times}A_4$ off-block); and the missing corner $F=R+Q$ has $\det B_F=52=\dim F_4$ with
$248=52+14+26\cdot7$ --- the genuine $E_8\!\supset\!F_4{\times}G_2$ branching $(52,1){\oplus}(1,14){\oplus}(26,7)$
as an operator \emph{shadow}. These are exact identities; their atlas/embedding significance is kept audit-level
(\mC), since the coefficients are themselves atoms.
\end{keybox}

\begin{keybox}[The Koide point is one branch of a double cover \mE\ (\veri{v81\_singular\_pencil\_matrices.py})]
The anchor-block determinant $\det B_{K+xQ}=(3x+2)(3x+5)$ is \emph{quadratic}, so
\[
  \boxed{\ y^2=\det B_{K+xQ}=(3x+2)(3x+5)\ }
\]
is a $2{:}1$ cover of the pencil line, \emph{ramified exactly over its two zeros}. Hence the Koide/gap point
$x=-|\Z_2|/\Nfam=-\tfrac23$ and the carrier point $x=-\gcar/\Nfam=-\tfrac53$ are \textbf{the two branch points of
one double cover}: \emph{Koide is literally ``the other side'' of the carrier point}. The deck involution is
$y\mapsto-y$ of degree $2=|\Z_2|$ --- the sheet itself --- which is the line-bundle reflection of the
$\mathrm{Spin}(10)\!\to\!\mathrm{SO}(10)$ half-spinor cover the carrier lives on. Three rigid readings:
\begin{itemize}\setlength\itemsep{1pt}
  \item \textbf{Branch locus $=$ spine ends / family.} The ramification is at $-2/3$ and $-5/3$, i.e.\
  $-(|\Z_2|,\gcar)/\Nfam$ --- the \emph{two ends} of the spine $(2,3,4,5)$ over $\Nfam$.
  \item \textbf{Discriminant $=\Nfam^4=81$} (a perfect square $\Rightarrow$ rational, split branch points);
  $\sqrt{81}=9=\Nfam^2=\det B_Q$.
  \item \textbf{Separation $=1=$ one transport period} (the $K\!\to\!L$ step $x:0\!\to\!1$); the deck-translation
  $x\mapsto x+1$ carries the carrier point onto the Koide point.
  \item \textbf{Branch divisor $=$ (scalaron, $\AL$).} Sum of branch points $=-7/3=-7/\Nfam$, product $=10/9=\AL/\Nfam^2$;
  the integer labels $\{|\Z_2|,\gcar\}=\{2,5\}$ have $\mathrm{sum}=7$, $\mathrm{prod}=10=\AL$, $\mathrm{diff}=3=\Nfam$.
  \textbf{So the scalaron is the \emph{trace} of the branch divisor} --- its several readouts ($\Oadm{-}10b_1$,
  $\AL{-}\Nfam$, $\gcar{+}|\Z_2|$, the mixed anchor area) collapse to this one geometric origin.
  \item \textbf{The sheet lives on the cut.} The fibre $y^2{=}\det B$ is $\AL$ at $K$ and $4\AL$ at $L$, so transport
  multiplies $y$ by $|\Z_2|$; the $\Z_2$ sheet endpoint $x{=}{-}1$ sits \emph{between} the branch points where
  $y^2{=}\det B_{K-Q}{=}{-}|\Z_2|{<}0$ --- on the imaginary branch cut joining Koide and the carrier.
\end{itemize}
\noindent At each branch point the pencil clears its $3$-denominator into an integer compiler matrix.
\emph{Koide} $C_{2/3}=3K-2Q$ reproduces the carrier glue: $\operatorname{tr}=15=\dim A_3$, $\textstyle\sum=45=\dim D_5$,
$\det=60=D_{\mathrm{start}}$, and the (now forced rank-one) anchor block
$B_{C_{2/3}}=\begin{psmallmatrix}45&55\\63&77\end{psmallmatrix}=\binom{\gcar}{7}\,(\,\Nfam^2\ \ 11\,)$ with
$\sum=240=|R(E_8)|$. \emph{Carrier} $C_{5/3}=3K-5Q$ is the charge-\emph{neutral} side ($\sum_{ij}=0$) with
$\chi(\lambda)=(\lambda+|\Z_2|)(\lambda^2+\lambda+|R^+(A_3)|)=(\lambda+2)(\lambda^2+\lambda+6)$, splitting off
the sheet eigenvalue $-|\Z_2|$. So the recurring $\tfrac23$ is not a quotient that ``appears everywhere'' --- it
is one ramification point of the mass$\leftrightarrow$transport double cover. The sharper open question is then
\emph{why the leptonic source-to-pole transfer renormalises onto a branch point, and which sheet ($y\to\pm$) the
parity selects} --- a cleaner target than ``why $2/3$''. (The $D_5/A_3$ and $E_8$ labels on the clearing matrices
are exact; their per-entry atlas reading stays audit-level \mC.) \emph{Gravitational twin:} the same split
$\Nfam$-linear double-cover structure reappears in classical gravity --- the Schwarzschild--de~Sitter mass
line has $\operatorname{disc}\propto(1-3m)(1+3m)$ (branch points $\pm1/\Nfam$, separation $2/3$, deck $=$
horizon swap), the Nariai roots are the traceless anchor $(1,1,-2)$, and $2/3$ is the exact Nariai/de-Sitter
entropy bound (Appendix~H; \veri{v101\_horizon\_anchor.py}).
\end{keybox}

\begin{keybox}[The Koide branch is a \emph{forced} attractor, and the clean split is non-generic \mE/\mC\ (\veri{v82\_koide\_attractor\_splitting.py})]
This sharpens the open question above. Work in the branch coordinate $q=-3x$ (Koide $q=2$, carrier $q=5$).
\textbf{(A) The source-to-pole map is forced, not postulated.} Any branch-divisor-preserving M\"obius map of the
pencil line that fixes \emph{both} branch points is, in the cross-ratio coordinate $\rho=(q-2)/(5-q)$, exactly
$\rho\mapsto\mu\rho$ --- hence \emph{unique} once $\mu$ is fixed. The multiplier is not a new input: $\mu=(2/3)^6$
is the subleading eigenvalue $\lambda_2$ of the \emph{already-established} unique gapped boundary transfer operator
(spectrum $\{1,(2/3)^6,(1/3)^6\}$, gap $6\log\tfrac32>0$; \veri{v54\_seam\_horizon\_keystones.py}, \veri{v56\_unique\_attractor.py}),
and the Koide branch $x=-2/3=-|\Z_2|/\Nfam$ is precisely its cusp weight $2/3=\{0,1,|\Z_2|\}/\Nfam$ (the carrier maps to
Koide under the deck period $q\mapsto q-\Nfam$, the same cusp class). With $\mu=(2/3)^6$ the unique map is the integer
M\"obius map
\[
  \boxed{\ F(q)=\frac{2(569q-3325)}{665q-3517}\ },\qquad F(2)=2,\ \ F(5)=5,\ \ F'(2)=(2/3)^6<1,\ \ F'(5)=(3/2)^6>1,
\]
so Koide is the \textbf{attractor} and the carrier the \textbf{repeller}, and $q_n\to2$ i.e.\ $x\to-2/3$ at the
\emph{same} rate that governs the SM flavor gap and horizon recovery. The three structural postulates needed for a
Koide derivation thus collapse to \emph{one} physical identification --- \emph{source-to-pole transfer $=$
branch-preserving relaxation of the established transfer operator} --- after which rate, uniqueness and direction
all follow. That single identification stays \mC; the rest is \mE. \textbf{(B) The clean rational split is
non-generic.} The splitting-type exponents $(2,1,1)$ (the anchor $a=(1,1,2)$) placed across the three family slots
give three different covers: $(1,1,2)$ gives $\det B_{K+xQ}=(3x+2)(3x+5)$, $\mathrm{disc}=81=\Nfam^4$ (\emph{split},
rational labels $\{2,5\}$); $(2,1,1)$ gives $(x+2)(9x+11)$, $\mathrm{disc}=49=7^2$ (\emph{split}); $(1,2,1)$ gives
$5x^2+10x+3$, $\mathrm{disc}=40=|R(D_5)|=\det B_L$, which is \emph{not} a square --- a \textbf{non-split} cover with
Galois-conjugate branch points $-1\pm\sqrt{10}/5$ centred exactly on the sheet $x=-1$ ($q=3$), separation
$2\sqrt{\AL}$. (Negative control, as for the cover itself: the $R$-pencil has $\mathrm{disc}=153$, also not a
square.) So Koide and the carrier appearing as \emph{separate rational} singularities is a special feature of the
canonical anchor placement, not a generic one --- a hardening of anchor-first. The physical fermion-sector
identification of the three placements is not claimed and stays \mC.
\end{keybox}
